\section{问题定义与业务时间对齐}

设预测日为 \(D\)，目标窗口为北京时间 \(D\) 日 00:00--23:00。本文直接预测区域级聚合出力序列 \(\mathbf{y}\in\mathbb{R}^{24}\)，即目标区域内 58 个风电场站在每个目标时刻的合计出力，不输出单场站功率。目标窗口前最近 16 步场站级功率作为状态上下文，记为 \(\mathbf{P}^{c}\in\mathbb{R}^{N\times16}\)，其中 \(N=58\) 为场站数量；模型输出为 \(\hat{\mathbf{y}}\in\mathbb{R}^{24}\)。

本文遵循业务化时间可用性约束。对于日样本评价，最新可用功率上下文覆盖 \(D-1\) 日 08:00--23:00。Pangu-Weather 预报在 \(D-1\) 日 08:00 起报，lead 0--15 对应目标窗口之前的后报气象段，lead 16--39 对应 \(D\) 日 00:00--23:00 的未来预报气象段。因此，气象输入包括后报气象 \(\mathbf{X}^{h}\in\mathbb{R}^{16\times C\times H\times W}\) 和未来气象 \(\mathbf{X}^{f}\in\mathbb{R}^{24\times C\times H\times W}\)，二者沿时间维拼接为 \(\mathbf{X}\in\mathbb{R}^{40\times C\times H\times W}\)。场站邻接矩阵 \(\mathbf{A}\in\mathbb{R}^{N\times N}\) 由近期场站功率相关性构造，用于功率状态编码。

训练阶段采用滑动窗口样本，以提高样本数量并覆盖更多预测起点；验证和测试阶段采用严格日样本，以模拟每天固定起报后的业务化预测流程。无论训练还是测试，模型均只使用目标窗口开始前真实可用的信息，避免未来功率泄漏。该设置使本文问题不同于离线随机切分的序列预测任务，也使功率状态条件化的作用能够在逐 horizon 分析中被直接检验。
